\section{Proof of the branch theorem}

Here $d_7=3$.  A cube in $\C(z)$ has divisor valuations divisible by three.
The exact valuations $\pm2$ supplied by \cref{prop:gcd} prove that $N_7$ is
not a cube.  Near a zero $\alpha$,
\[
 G_7(z)=(z-\alpha)^{2/3}h(z),
 \qquad h(\alpha)\ne0,
\]
and a pole gives order $-2/3$.  A single-valued meromorphic function has
integral local order.  This proves \cref{thm:main}.

The obstruction is stronger than the statement that a convenient rational
formula was not found.  It is invariant under rational re-expression of the
norm and is already present at the first split prime.  Any ordinary
finite-dimensional determinant ratio has integral local orders and therefore
cannot equal $G_7$.

It would be incorrect, however, to infer a global natural boundary.  The
origin branch remains holomorphic on $|z|<1$, and the C45 prime product remains
holomorphic on $\operatorname{Re}s>1/2$.  Cross-prime monodromy, global
cancellation, and alternative operator categories have not been classified.

